\chapter{纯净物其实很难“纯”}

\begin{kaichang}
走进超市的饮料货架，你至少能找到三种“水”：矿泉水，瓶身上印着钙、镁、钾的含量表；饮用纯净水，包装上写着“层层净化”；药店和五金店里还卖蒸馏水，买它的人多半是为了加湿器、汽车电瓶或者实验室。

先别往下读，猜一猜：这三种水，哪一种最“纯”？如果买到一瓶真正合格的蒸馏水，它是不是就“只有水，别的什么都没有”？“纯”这件事，有没有一个能做到头的终点？

把答案写在纸上。这一章结束时我们回来对答案——你会发现，“纯净”是化学里最常用、也最容易被想歪的概念之一。
\end{kaichang}

\section{看上去“一样”，不算数}

先回忆第 1 章分好的两类变化。白糖溶进水里，糖没有变成新物质，只是“藏”进了水；糖加热变焦，才是变成了别的东西。前一类的产物有个共同点：它是两种（或更多）东西\textbf{混在一起}，谁也没变成谁。化学给这种状态起了名字：

\begin{itemize}[itemsep=2pt]
\item \textbf{纯净物}：只由一种物质组成的东西。蒸馏水接近它，氧气、铜丝、精盐也接近它。
\item \textbf{混合物}：由两种或更多种物质混在一起组成，各成分保持自己本来的面目。空气、海水、牛奶、你手里的不锈钢勺子，都是混合物。
\end{itemize}

听起来简单，但“怎么判断一样东西是纯的还是混的”并不简单。眼睛是最靠不住的裁判。一杯白糖水清澈透明，看上去“只有一种东西”，可你蒸干水分，锅底会留下白糖——它一直是两户人家。牛奶看上去也很均匀，但把它静置几天会分层，用光一照又白又浑（光被里面的小颗粒散射得到处都是），跟清水完全两样。

所以侦探要靠\textbf{行为}而不是\textbf{长相}来断案：静置会不会分层？能不能被滤纸拦住？蒸干有没有残渣？加热时是不是在一个固定温度沸腾？这些可观察、可测量的线索，才是“纯不纯”的判决依据——这正是第 2 章说的“用间接证据推理看不见的东西”。

混合物内部还能再分两档。像盐水、空气、均匀的茶水这样，成分分布得处处一样的，叫\textbf{均一混合物}（也叫溶液，不只限于液体）；像泥沙水、油醋汁、花岗岩这样，一眼能看出这一坨那一坨的，叫\textbf{不均一混合物}。但要注意：均一不等于纯净。盐水再均匀，里面也是水和盐两户人家。

现在可以给开头点名的四位“嫌疑人”定性了：

\begin{itemize}[itemsep=2pt]
\item \textbf{空气}：均一的气体混合物，主要住户是氮气和氧气，还有氩气、二氧化碳和水蒸气。
\item \textbf{牛奶}：不均一的混合物，是水、脂肪小球、蛋白质、乳糖等的大杂烩，均匀程度介于中间地带。
\item \textbf{海水}：均一混合物，一碗汤里溶解着几十种盐，主要是氯化钠。
\item \textbf{合金}（如不锈钢、黄铜）：均一混合物，几种金属（有时掺非金属）熔在一起，均匀到你看不出任何接缝。
\end{itemize}

还有一条很实用的现象判据：\textbf{纯净物有固定的熔点和沸点，混合物没有}。纯冰在 $0\,^{\circ}\mathrm{C}$ 整开始融化，熔化过程中温度一动不动；而一块黄油（油脂混合物）是在一段温度区间里慢慢软化的。纯水在 $100\,^{\circ}\mathrm{C}$（标准大气压下）沸腾，糖水却是越煮沸点越高——因为水不断跑掉，剩下的糖水越来越浓。冬天往路面撒盐化冰，利用的正是“掺了杂质熔点下降”这条混合物的脾气。状态变化背后的粒子道理，第 4 章马上讲。

\section{把镜头推进到粒子}

宏观分类只是现象层的说法。换上第 2 章起就在用的粒子眼镜，定义立刻变得锋利起来：

\textbf{纯净物，就是只含同一种粒子的东西。}一杯纯水（理想情况下）从头到尾都是水分子，翻来覆去就那一种；一根铜丝从头到尾都是铜原子。\textbf{混合物，就是两种以上不同的粒子混在一起，各自还是各自}：空气里是氮分子、氧分子、氩原子、二氧化碳分子在并肩乱撞；海水里是水分子中间游荡着各种盐的带电粒子（离子）。

纯净物里还能再分一层，标准是看这种粒子里有\textbf{几种原子}：

\begin{itemize}[itemsep=2pt]
\item \textbf{单质}：粒子只由同一种原子组成。氧气分子由两个氧原子搭成，铜丝里是铜原子的方阵，金刚石里全是碳原子。
\item \textbf{化合物}：粒子由不同种原子按固定比例牢牢结合成一个整体。一个水分子是两个氢原子加一个氧原子；一粒食盐晶体是钠和氯按一比一交替排成的格子（第 18 章会细看这种格子）。水的氢和氧拆不开——除非发生化学变化，第 1 章的界线在这里又一次出现：分开混合物靠“挑拣”，拆开化合物必须“动手术”。
\end{itemize}

用粒子语言重述四位嫌疑人：牛奶里是巨大的脂肪小球和蛋白质长链漂浮在水分子的海洋里，颗粒比分子大成千上万倍，所以它会散射光、会分层——它处在“均一”和“不均一”之间的灰色地带，化学上叫胶体。不锈钢里是铁、铬、镍的原子直接混坐在同一片晶体方阵里，均匀到原子尺度。这两例提示我们：均一不均一的分界，本质上是“混合的颗粒有多大”——从几个分子到肉眼可见的小液滴，是一条连续的谱，不是两格抽屉。

还要强调一句最容易被忽略的话：\textbf{混合的时候，粒子本身没有变}。氮分子混在氧分子中间，还是原来那个氮分子；盐溶进水里，钠离子也还是那个钠离子。混合只是“住在一起”，不是“结婚”。这就是为什么混合物没有固定的新性质——盐水咸、水不咸，各管各的账。

\section{分离：利用差异，而不是变魔术}

既然是“住在一起没结婚”，那想把混在一起的几户人家分开，就不需要打碎任何东西，只需要找到它们之间\textbf{某一个性质上的差异}，再设计一个装置把这个差异放大成“你走这边、它走那边”。这是本章最重要的一条规则：\textbf{一切分离方法，都是对某种性质差异的利用}。

\begin{itemize}[itemsep=2pt]
\item \textbf{过滤}：利用\textbf{颗粒大小}的差异。滤纸的小孔放行水分子和溶解的小粒子，拦住大颗粒。咖啡滤纸拦咖啡渣、豆浆机筛豆渣、口罩拦飞沫，都是同一招。
\item \textbf{蒸馏}：利用\textbf{沸点}的差异。加热时沸点低的成分先变成气体跑掉，再把气体冷却收回液体。海水淡化、酒厂蒸出酒精、实验室制蒸馏水，全靠它。
\item \textbf{萃取}：利用\textbf{“更爱待在哪种溶剂里”}的差异。同一个溶质，在水里和在油里溶解能力天差地别。用酒精浸泡花瓣提取香气和色素、用适当的溶剂从植物里提取药物成分（青蒿素的提取就用过类似思路），都是让目标分子“搬家”到更爱它的那一边。
\item \textbf{层析}（色谱）：利用\textbf{“恋家程度”}的微小差异。让混合物随流动的液体（或气体）经过一个固定表面，爱粘附在表面上的成分走得慢，爱跟着流体跑的成分走得快，几户人家就沿着一条带子拉开了距离。墨水的颜色分离就是层析。
\item \textbf{结晶}：利用\textbf{溶解度随温度（或水量）变化}的差异。盐田里晒海水，水蒸发了，盐没地方待，只好排队析出成晶体；热浓糖水冷下来会结出糖晶，是同一原理。
\item \textbf{离心}：利用\textbf{密度和沉降速度}的差异。机器高速旋转，重的成分被甩到外圈、沉到底部。医院抽一管血离心后分成血浆和血细胞两层，用的就是它。
\end{itemize}

把这些整理成一张“侦探工具表”：

\begin{table}[htbp]
\centering
\begin{tabular}{lll}
\toprule
方法 & 利用的差异 & 生活或工业例子 \\
\midrule
过滤 & 颗粒大小 & 咖啡滤纸、净水滤芯 \\
蒸馏 & 沸点高低 & 海水淡化、蒸馏水 \\
萃取 & 在不同溶剂中的溶解度 & 提取香料、药物成分 \\
层析 & 对固定表面的吸附强弱 & 分离墨水、药检 \\
结晶 & 溶解度随温度、水量的变化 & 盐田晒盐、冰糖 \\
离心 & 密度与沉降快慢 & 血液分层、甩干机 \\
\bottomrule
\end{tabular}
\end{table}

最后守一条底线：\textbf{分离不创造物质，也不消灭物质}。盐水分离成盐和纯水，盐粒一颗不多、一颗不少，水分子一个没丢；所有原子只是被“重新分了房间”。这条守恒的账本思想现在先记住，到第 26 章它会成为化学方程式记账的绝对底线。

\begin{qiaoliang}
“只剩一点点杂质”到底是多少？来算一笔账。一瓶饮用纯净水约 \ud{500}{g}，合格的纯净水要求杂质（总溶解固体）低于约 $10$ ppm——ppm 是“百万分之一”的意思，即每一百万份质量里最多 10 份杂质。就算按 $10$ ppm 算，这瓶水里杂质的质量也只有：
\[500\ \mathrm{g} \times \frac{10}{10^{6}} = 5\times 10^{-3}\ \mathrm{g}\]
五毫克，肉眼根本看不见。听起来“纯到家了”？换到粒子世界再算。一个水分子的质量约 $3\times 10^{-23}$ g，所以 $500$ g 水约有
\[\frac{500}{3\times 10^{-23}} \approx 1.7\times 10^{25}\ \text{个水分子}\]
杂质粒子的质量与水分子差不了几个数量级，于是那五毫克杂质约含 $10^{20}$ 个粒子——一亿亿个的一百倍。结论有点反直觉：宏观上“几乎为零”的微量，在粒子尺度上依然是人山人海。“微量”和“没有”之间隔着天文数字，这就是为什么检测痕量污染物那么难、也那么重要。这种“先估数量级再下结论”的算法，第 5 章会配上正式工具重新操练。
\end{qiaoliang}

\begin{shiyan}
\textbf{动手试一试：把黑色墨水“拆开”——纸上层析}

材料：黑色（或深色）水性笔一支、滤纸一张（或吸水性好的厨房纸巾）、玻璃杯一个、清水、铅笔、剪刀、一枚回形针或一根牙签。

步骤：
\begin{enumerate}[itemsep=1pt]
\item 把纸裁成宽约 2 cm、长约 10 cm 的纸条，用铅笔在距一端约 2 cm 处轻轻画一条横线（铅笔的痕迹不会被水带着跑）。
\item 用黑笔在横线中央点一个浓而小的墨点，晾干；想效果明显，可以干了再点一次。
\item 往杯里倒约 1 cm 深的清水。把纸条没点墨的一端用牙签架在杯口，让画线一端垂入水中——\textbf{水面必须低于铅笔线和墨点}。
\item 静置观察 10～20 分钟。水会顺着纸向上爬，经过墨点时把色素带着一起走。
\end{enumerate}

观察点：墨点会沿纸条拖出几条颜色不同的色带——黑色墨水多半会分出蓝、紫、黄等颜色，而且各颜色爬的高度不一样。爬得快的颜色更“爱”跟着水跑，爬得慢的更“恋”纸面。一支笔里竟住着好几种染料，眼睛从来没告诉你。

安全提示：只用水和普通水性笔，本实验没有危险试剂；不要把纸和墨点直接按进水里（墨点会溶掉，实验失败）；做完洗手即可。若老师演示用酒精等溶剂的高级版本，请在实验室看老师操作，不要在家尝试。
\end{shiyan}

\section{现在，请符号登场}

概念站稳了，化学的速记符号才好出场。符号层的规矩是：\textbf{只有纯净物才有自己的化学式}，因为化学式写的是“那一种粒子长什么样”。

水分子是一个氧原子牵着两个氢原子，写作 \chem{H_2O}；氧气分子是两个氧原子，写作 \chem{O_2}；二氧化碳是一个碳原子加两个氧原子，写作 \chem{CO_2}。这些都是化合物——粒子里不止一种原子。氧气和铜（\chem{Cu}）、金刚石（\chem{C}）则是单质——粒子里只有一种原子。

食盐写作 \chem{NaCl}，但要小心：食盐里并不存在一颗颗“氯化钠分子”，而是钠离子和氯离子按一比一交替堆成的晶体格子，\chem{NaCl} 记录的是这个固定的比例（第 18 章会细讲）。

混合物呢？\textbf{写不出一个化学式}。空气没有化学式，你只能列成分表：按体积算，氮气 \chem{N_2} 约占 $78\%$，氧气 \chem{O_2} 约占 $21\%$，氩 \chem{Ar} 约占 $0.9\%$，二氧化碳 \chem{CO_2} 约占 $0.04\%$，外加多少不等的水蒸气。注意这个列表本身就在爆料：连“纯净的空气”都是四五种气体的混合物——那到底还有没有真正纯净的东西？这正是下一个故事的主角。

\section{小数点后第三位里藏着一种新元素}

\begin{zhengju}
19 世纪末，化学家们自信空气已经被摸清了：除掉氧气和水蒸气，剩下的就是“纯氮气”。英国物理学家瑞利勋爵（Lord Rayleigh）在 1892 年做一件看似枯燥的事——精确称量气体的密度，却撞见了麻烦。他从空气里除去氧气、二氧化碳和水蒸气后得到的“氮气”，密度是 \ud{1.2572}{g/L}；而用化学方法（分解氨）制出的氮气，密度是 \ud{1.2508}{g/L}。两者只差 $0.5\%$，小数点后第三位才看得出分歧。

换个马虎的人，大手一挥：“实验误差嘛。”瑞利偏不。他先把所有平庸的解释挨个排除：是不是仪器没校准？他换装置反复测，差异稳定存在。是不是空气氮里残留了氧气或水？他设计实验逐一清除，差异还在。是不是化学氮里混了更轻的氢气？他换了好几种化学来源制氮，结果始终更轻。排除到自己再无借口之后，他做了一个很“科学家”的动作：在《自然》杂志上公开登出这个矛盾，征求解释——他没有藏着这个“丢人的误差”。

化学家拉姆齐（William Ramsay）接住了难题。他提出一个当时听起来荒唐的替代解释：\textbf{空气里的“氮气”根本不纯}，里面混着一种比氮重、从不和任何物质反应的未知气体。检验办法很直接：让空气氮反复通过灼热的镁，氮会被镁“吃掉”（生成氮化镁），如果真有不为所动的东西，它应该剩下。结果，约有原来体积 $1/80$ 的气体怎么也不肯反应。拉姆齐用光谱一照——这种气体发出的彩色谱线和所有已知元素都对不上号。1894 年，两人宣布发现新元素氩：名字来自希腊语“懒惰”，因为它什么都不肯反应。

这个故事值得记住三点。第一，“纯氮”这个共识错了一百多年，错的幅度只有百分之一——\textbf{你宣称的纯度，永远不会超过你的测量精度}。第二，异常数据不是垃圾，是线索；瑞利若把“误差”抹平，氩还要再等很多年。第三，科学共同体的纠错靠的是可重复的数字，不是权威：任何人只要测得同样准，都能看到那个 $0.5\%$。而氩的“懒惰”性格，要到第 10 章的元素街区里才能找到根源。
\end{zhengju}

\section{“纯”的边界：这是一个连续谱}

现在可以给“纯净”画边界了。本章的模型——把物质分成纯净物和混合物两格——在大多数情况下好用，但它有三个必须知道的边界。

\textbf{边界一：纯度是连续的刻度，不是非黑即白的标签。}试剂瓶上的等级就是证据：分析纯约 $99.5\%$ 以上，优级纯更严，半导体工业要的硅则纯到 $99.999999999\%$——十一个 9，意思是每一千亿个硅原子里杂质原子不超过一个。为什么要这么疯？因为一个晶体管只有几百个原子宽，杂质原子会故意或意外地改变电流的走法（第 16 章和后面的材料章节还会遇到它）。所以“纯不纯”从来不是哲学问题，而是工程问题：\textbf{够不够用，由用途决定}。

\textbf{边界二：连“同一种分子”都还有暗层——同位素。}纯水里的所有分子确实都是水，但氢原子有两副面孔：绝大多数是普通氢，约每 3200 个氢原子里有 1 个是带一个中子的“重氢”（氘）。于是纯水里天然混着极少量“重水分子”。这算不纯吗？化学上通常不算——同位素的化学脾气几乎一模一样（第 8 章细讲它们的差别和用处）。但你看，连“百分之百的水”这个句子都经不起显微镜式的追问。

\textbf{边界三：均一和不均一之间没有刀口。}牛奶、豆浆、雾气、血液这样的胶体，颗粒比分子大得多、又比泥沙小得多，静置很久也不分家。它们算均一还是不均一？答案取决于你拿什么尺去量：肉眼看着均一，显微镜下就露馅。两格抽屉装不下连续的世界，这是所有分类模型的通病——好用，但别把它当成世界本身的法律。

\begin{wuqu}
“越纯越好、越纯越安全”“天然提取的总比化学合成的纯净健康。”——两个流传极广的误解。

先看反例：纯氧听起来“干净又高级”，但它会让本来看似温和的东西猛烈燃烧，医院里氧气瓶的管理比普通药品严得多；乌头碱、尼古丁都是“天然植物提取”且可以很纯，毒性却极强。反过来，合成的抗坏血酸（维生素 C）和从柠檬里提取的，是同一种分子，身体根本分不出“出身”。

再纠正一个方向性错误：\textbf{安全性取决于物质本身和剂量，不取决于“纯”或“天然”的标签}。“纯天然”不等于纯净——天然精油动辄含几百种分子，远比一瓶分析纯试剂“杂”；“化学成分”也不等于有害——水、盐、糖都是化学成分。至于“蒸馏水最纯所以最养人”，日常饮水里那点矿物质主要靠吃饭就能补齐，喝纯净水没有问题，但“越纯越健康”的推理在科学上是不成立的。
\end{wuqu}

\section{回到货架上的三瓶水}

可以兑现开场的承诺了。

\textbf{矿泉水}是一种均一混合物：水分子中间溶解着钙、镁、钾等矿物离子，瓶身上的含量表就是成分清单——它的卖点恰恰是自己“不纯”。\textbf{饮用纯净水}去掉了绝大部分离子和微生物，但仍溶解着空气里的氧气和二氧化碳，可能还有生产设备留下的痕量物质，纯度由检测标准界定，而不是“绝对只有水”。\textbf{蒸馏水}最接近理想中的纯净物：蒸馏把不挥发的盐类留在了锅里。但它依然不是终点——开盖后会溶入空气，容器会析出痕量成分，而且别忘了同位素：再怎么蒸馏，每三千来个水分子里照样有一个“重水分子”在滥竽充数。

所以开场问题的答案是：蒸馏水最纯，但“只有水，别的什么都没有”的状态不存在。“纯”不是一枚可以达到的奖牌，而是一把带刻度的尺——科学问的不是“纯不纯”，而是“纯到什么程度、用什么方法测的、对这个用途够不够”。换成这五个贯穿问题来说，本章其实只回答了一个开头：它由什么组成？而“组成部分怎么连接、怎么排列”，正是后面二十多章的主题。

\section{提纯是一门大生意，粒子运动是幕后老板}

别以为“纯”只是实验室的洁癖。你手里的手机是提纯工业的胜利：芯片用的硅从石英砂起步，先还原成纯度约 $98\%$ 的粗硅，再经过一系列化学转化和蒸馏（又是蒸馏！）提纯到十一个 9，最后还要用一种叫“区域熔炼”的办法，让熔融区沿硅棒缓慢扫过，把残余杂质像扫地一样赶到末端切掉。医学上，层析早已是侦探标配：药厂用它检查每一批药的纯度和杂质谱，反兴奋剂实验室用它从运动员的尿样里捞出十亿分之一克级别的违禁分子——还记得吗，“微量”在粒子世界也是人山人海，所以这点量足够定罪。环境监测站测自来水，同样是先分离、后定量。

而这些分离方法凭什么奏效？蒸馏靠沸点，沸点由粒子挣脱彼此的能力决定；结晶靠溶解度随温度变化，温度本质是粒子运动的剧烈程度；离心靠沉降，沉降是密度和运动的赛跑；层析靠“走还是留”的概率。所有线索都指向同一个幕后老板：\textbf{粒子永不停歇的运动}。下一章——第 4 章，我们就去直面它：为什么同一种物质有时是固体、有时是液体、有时是气体？状态的变化背后，粒子到底在干什么？至于溶解这件事本身的机关，第 25 章会专门拆开来讲。

\begin{tiaozhan}
层析凭什么能分开性质几乎一样的两种分子？秘密在“重复放大”。想象纸面由 $n$ 级台阶组成，分子在每级台阶上要么跟着水前进一步，要么被纸面吸附住停一停。设甲分子每级“前进”的时间比例是 $0.50$，乙是 $0.45$——单看一级，差距只有 $10\%$，根本分不开。但 $n$ 级之后，两种分子平均走到的位置正比于 $0.50n$ 和 $0.45n$，差距随 $n$ 线性拉大；而每条色带本身的宽度只随 $\sqrt{n}$ 变宽（随机涨落的规律）。于是：分离的差距按 $n$ 增长，模糊的程度只按 $\sqrt{n}$ 增长，台阶足够多级时，再像的两种分子也会被干净地撕成两条带。工业色谱柱的“台阶”可以有几万级，这就是它恐怖分辨能力的来源。同样的“微小差异乘以巨大次数”逻辑，到第 29 章讲反应快慢、后面讲进化时还会回来。跳过本框不影响主线。
\end{tiaozhan}

\section*{练一练}

\begin{sikaoti}
\item 给下列各物分类（纯净物/混合物；纯净物再分单质/化合物），并各画一幅“粒子示意图”：洁净的空气、蔗糖晶体、不锈钢、牛奶、氧气。在图旁注明：圆圈大小不代表原子的真实比例。
\item 一杯白糖水和一杯汽水都是均一混合物。画两幅粒子图，各标出至少三种不同的粒子，并用一句话回答：为什么“均一”不等于“纯净”？
\item 你拿到一把混有沙子和铁粉的食盐，想回收食盐。设计分离流程，画成物质流图（每一步注明利用了什么性质差异），并说明：整个过程中原子的总数发生了什么变化？
\item 纸上层析中，蓝色带爬得比黄色带高。哪一种色素更“恋”流动的水、哪一种更“恋”纸面？如果换用长一倍的纸，预测两条色带的间距会怎样变化，并说出理由（可借用挑战框的思路）。
\item 假如瑞利的仪器只能测到小数点后一位，氩还会被发现吗？用本题说明“测量精度”和“纯度宣称”之间的关系。
\item 有广告说某产品“$100\%$ 纯净，绝无任何化学物质”。请从两个角度反驳：一是“纯净”为什么做不到百分之百（举同位素或溶解气体的例子）；二是“绝无任何化学物质”这句话本身错在哪里。
\end{sikaoti}
